\documentclass[instructions.tex]{subfiles}
\begin{document}
 
\chapter{La confusion des pécheurs au jour du jugement et leur désespoir}
 
La vue des péchés manifestés à tout l’univers, accablera de confusion les pécheurs; la sentence prononcée les jettera dans le plus affreux désespoir.
 
\primo  Les rayons du soleil ne pénètrent pas si clairement un cristal, que les yeux de Dieu pénétreront le fond des cœurs; c’est pour le jour du jugement qu’il s’en est réservé la manifestation. En ce jour de révélation, Jésus-Christ manifestera jusqu’aux péchés les plus secrets, avec leurs circonstances les plus honteuses.
 
On verra alors ce qu’il y a de plus caché; tout ce qui s’est passé dans les ténèbres et dans les lieux écartés; tout ce qu’on a fait, ce qu’on a dit, ce qu’on a pensé dans les compagnies, dans les rendez-vous, dans les promenades et les jeux. On verra les désirs de plaire, les intrigues, les attaches et les crimes honteux sur lesquels on s’aveuglait. On verra les projets de chicane et les intentions obliques dans les entreprises, les vols secrets et les fourberies, les artifices de la cupidité et de l’intérêt, pour s’enrichir aux dépens d’autrui; les impostures et les jalousies cachées au fond du cœur; l’hypocrisie enfin de tant de personnes qui cachent leurs vices sous l’apparence de l’honnête homme, tout sera manifesté.
 
Un témoin qui vous surprendrait dans le crime, vous ferait rougir. Quelle sera votre confusion, vous qui craignez les discours du monde, qui craignez même de déclarer vos faiblesses à un confesseur plein de charité, lorsque les anges, les saints, tous les hommes, connaîtront ce qu’il y a de plus horrible et de plus honteux dans votre conduite!
 
Les pécheurs regarderaient l’enfer comme un lieu de repos, dans ce jour épouvantable, s’ils avaient la liberté de demeurer cachés; mais il leur sera autant impossible de se cacher au jugement, qu’il leur sera insupportable d’y paraître. Ils seront forcés d’y comparaître comme des criminels pour se voir confondus devant l’univers par un Dieu tout-puissant. Oh! quelle consternation, quand ils se verront accablés des reproches de ce juge terrible, des reproches des anges et des saints, des reproches des démons! Que reviendra-t-il alors du péché, que la honte et l’horreur de l’avoir commis?
 
\secundo La sentence de réprobation mettra le comble au désespoir de ces misérables. Jésus-Christ, d’une parole, renversa les soldats qui venaient le saisir au jardin des Olives. Ah! si la majesté de cet Homme-Dieu fut si redoutable dans un temps où il paraissait comme Sauveur, que sera-ce lorsqu’il parlera en juge, et qu’il prononcera, d’une voix plus effrayante que les tonnerres, ces terribles paroles : Retirez-vous de moi, maudits, allez au feu éternel\footnote{Discedite à me, maledicti, in ignem æternum. Matth. 25. 41.}.
 
L’excès de leur désespoir sera de voir que celui qui les condamne est leur rédempteur, qui leur mettra devant les yeux les grâces, les moyens de salut qu’il a donnés, le sang qu’il a répandu, les tourmens qu’il a soufferts pour les sauver. Oh! l’affligeante situation, dit saint Eucher, de voir son Sauveur et de le perdre en même temps, et de périr à la vue d’un Dieu qui voulait nous sauver\footnote{Quàm lugubre erit homini Deum videre et perdere, et antè pretium sui perire conspectum. Euch.}.
 
S’il est si affligeant de se voir séparé d’une personne qu’on aime, que sera-ce d’être pour toujours séparé de son Dieu, le centre de notre âme et son souverain bien! Séparation cruelle et désespérante, qui sera le plus insupportable tourment de ces misérables. Mille carreaux et mille éclats de foudre, lancés sur une même tête, dit saint Chrysostome, ne sont pas si redoutables qu’un seul regard de cette face divine, qui se détournera du pécheur avec indignation.
 
Quel désespoir de se voir rejeté pour jamais de cette présence adorable, d’avoir tout perdu en perdant la possession de son Dieu, et la société des saints, pour s’être livré à des vanités, à des plaisirs, à des amusemens qui ne sont plus, qui n’avoient qu’une apparence passagère et trompeuse! Faites par votre grâce, ô mon Dieu! que je ne vous perde pas en cette vie par le péché, afin que je ne sois pas alors séparé de vous. « O Jésus! souvenez-vous du sang que vous avez répandu pour moi; ne me perdez pas en ce jour terrible\footnote{Recordare, Jesu pie, quòd sum causa tuæ viæ; ne me perdas illà die. Prose.}. »
 
\section{Résolutions}
 
\primo Puisque, selon la parole de Jésus-Christ, rien ne demeurera caché au jugement général, il faut que je ne me permette désormais plus rien, ni dans mes pensées, ni dans mes affections, ni dans mes discours, mes regards et toutes mes démarches, qui puisse me causer de la confusion.
 
\secundo Et afin d’éviter la sentence de réprobation, il faut que je répare le passé par une confession entière, accompagnée de sentimens de douleur qui ne finiront qu’avec ma vie.
 
\end{document}
